\chapter*{ABSTRAK}
\addcontentsline{toc}{chapter}{ABSTRAK}

\textit{Penelitian ini merancang sebuah prototipe sistem verifikasi dokumen akademik berbasis
blockchain untuk mengatasi kelemahan sistem terpusat yang rentan terhadap Single
Point of Failure (SPoF) dan manipulasi data sistem. Sistem menggunakan privat
blockchain Hyperledger Besu dengan konsensus QBFT serta protokol data terstruktur
EIP-712 untuk menghasilkan tanda tangan digital yang human-readable. Data yang
disimpan di blockhain direpresentaikan sebagai identityHash dan fileHash, kemudian
ditandatangani dan disimpan di blockchain melalui API Middleware. Selain itu sistem
menerapkan mekanisme zero-gas dan pemisahan antara relayer dan signer. Hasil
pengujian menunjukan bahwa sistem memenuhi parameter transaksi blockchain
berdasarkan NISTIR 8202, serta mampu memverifikasi integritas dokumen melalui
pencocokan idendityHash dan fileHash dengaan data on-chain. Selain itu, transaksi juga
berhasil dilakukan tanpa ada biaya gas pada jaringaan. Kode implementasi bisa diakses
di halaman: https://github.com/rizkycahyono97/tasdiqi.}\\	

\noindent\textit{\textbf{Keywords:} Blockchain, Hyperledger Besu, EIP-712, Verifikasi Dokumen Akademik, Tanda Tangan Digital}\\	%edited


